\documentclass[11pt]{article}
\usepackage[a4paper,margin=1in]{geometry}
\usepackage{amsmath,amssymb,amsthm,mathtools}

\newtheorem{theorem}{Theorem}[section]
\newtheorem{lemma}[theorem]{Lemma}
\newtheorem{proposition}[theorem]{Proposition}
\newtheorem{conjecture}[theorem]{Conjecture}
\newtheorem{remark}[theorem]{Remark}

\newcommand{\R}{\mathbb R}
\newcommand{\Om}{\Omega}
\newcommand{\dT}{\delta_T}
\newcommand{\A}{\mathcal A}
\newcommand{\capr}{\operatorname{cap}}
\newcommand{\HH}{\mathcal H}

\title{Planar Radialization Lemma: Proof Attempt}
\author{}
\date{}

\begin{document}
\maketitle

\section{Purpose}

The hybrid route would be closed by the following statement.

\begin{conjecture}[Planar radialization at torsion scale]
\label{conj:rad}
Let \(E=E_{\hat t}\subset\R^2\) be the regular good level obtained from
the corrected Step 1, normalized by \(|E|=\pi\), with
\[
   |\Omega\setminus E|\le C\dT(\Omega),
   \qquad
   \dT(E)\le C\dT(\Omega).
\]
Then there exists a small-amplitude radial graph \(G\), after a
translation and area normalization, such that
\[
   |E\Delta G|\le C\dT(\Omega),
   \qquad
   \dT(G)\le C\dT(\Omega).
   \tag{R}
\]
\end{conjecture}

If true, Conjecture~\ref{conj:rad} closes the \(n=2\) hybrid route:
the existing small-amplitude radial-graph theorem gives
\(\A(G)^2\le C\dT(G)\), the symmetric-difference estimate transfers this
to \(E\), and Step 5 transfers from \(E\) to \(\Omega\).

This note tries to prove Conjecture~\ref{conj:rad}.  The outcome is
partial: two needed analytic pieces are proved below, and the proof is
reduced to two explicit geometric lemmas.  I do not see a way to remove
those lemmas using only the current repo.

\section{A capacity identity that is genuinely useful}

Let
\[
   T(D):=\int_D u_D\,dx=\int_D|\nabla u_D|^2\,dx
\]
be the torsional rigidity, where \(-\Delta u_D=1\) in \(D\) and
\(u_D\in H^1_0(D)\).

\begin{lemma}[Exact torsion loss under an internal obstacle]
\label{lem:caploss}
Let \(D\subset\R^2\) be open with finite measure and let
\(K\Subset D\) be compact.  Put \(D_K:=D\setminus K\).  Then
\[
   T(D)-T(D_K)
      =\int_D |\nabla(u_D-u_{D_K})|^2\,dx .
   \tag{2.1}
\]
Consequently, if \(u_D\ge a>0\) quasi-everywhere on \(K\), then
\[
   T(D)-T(D_K)\ge a^2\,\capr_D(K),
   \tag{2.2}
\]
where \(\capr_D(K)\) is the relative condenser capacity of \(K\) in \(D\).
\end{lemma}

\begin{proof}
Let
\[
   J_D(v)=\frac12\int_D|\nabla v|^2\,dx-\int_D v\,dx .
\]
The minimizer of \(J_D\) over \(H^1_0(D)\) is \(u_D\), and
\(J_D(u_D)=-T(D)/2\).  Since \(u_{D_K}\in H^1_0(D)\), and since
\(J_D(u_{D_K})=J_{D_K}(u_{D_K})=-T(D_K)/2\), strict convexity around
the minimizer gives
\[
   J_D(u_{D_K})-J_D(u_D)
      =\frac12\int_D|\nabla(u_D-u_{D_K})|^2\,dx .
\]
This proves (2.1).

For (2.2), \(h:=u_D-u_{D_K}\) is admissible for the condenser problem
after division by \(a\): it is in \(H^1_0(D)\) and satisfies
\(h\ge a\) quasi-everywhere on \(K\).  Hence
\[
   \int_D|\nabla h|^2\,dx\ge a^2\capr_D(K).
\]
\end{proof}

This is the right mechanism for inward fjords.  If a missing piece \(K\)
sits a fixed positive torsion height inside the filled set, then
removing it has a capacity cost.  In the plane, capacity is much larger
than area for tiny holes, so a small torsion deficit forces the area
that must be filled to be small.

\section{Filling is compatible with the torsion deficit}

The next observation removes a previous worry: if the repair is mostly
``fill holes'', then the torsion deficit of the repaired set remains
controlled after rescaling.

\begin{lemma}[Monotone filling and area rescaling]
\label{lem:fillscale}
Let \(E\subset H\subset B_R\) be open planar sets with
\[
   |E|=\pi,\qquad |H|=\pi+v,\qquad 0\le v\le v_0(R).
\]
Let \(s=(\pi/|H|)^{1/2}\) and \(G=sH\), so \(|G|=\pi\).  Then
\[
   \dT(G)\le C_R\bigl(\dT(E)+v\bigr),
   \tag{3.1}
\]
where \(\dT(D)=T(B_{|D|})-T(D)\), up to the fixed harmless normalization
used in the repo.
\end{lemma}

\begin{proof}
By monotonicity, \(T(H)\ge T(E)\).  Also
\[
   T(B_{|H|})-T(E)
      =\bigl(T(B_{|H|})-T(B_\pi)\bigr)+\bigl(T(B_\pi)-T(E)\bigr)
      \le C v+\dT(E),
\]
because in dimension two \(T(B_r)=\pi r^4/8\), while
\(|B_r|=\pi r^2\), so \(T(B_{|H|})\) is \(C^1\) in the area near
\(\pi\).  Therefore
\[
   T(B_{|H|})-T(H)\le C v+\dT(E).
\]
Scaling by \(s\) multiplies torsional rigidity by \(s^4\), and it sends
the equal-volume ball for \(H\) to \(B_\pi\).  Thus
\[
   \dT(G)=s^4\bigl(T(B_{|H|})-T(H)\bigr)
      \le C_R(\dT(E)+v).
\]
\end{proof}

Thus a proof of Conjecture~\ref{conj:rad} can focus on producing a
small-amplitude radial graph hull \(H\supset E'\), where \(E'\) is \(E\)
after deleting only low-torsion outward debris, with total fill/delete
volume \(O(\dT(E))\).

\section{What would close the proof}

The previous lemmas show that the radialization proof would follow from
two geometric statements.

\begin{conjecture}[Planar core/capacity fill lemma]
\label{conj:corefill}
Let \(E=E_{\hat t}\) be the Step-1 good level in dimension two, with
\(\dT(E)\le\delta\) and \(D_H(\hat t)+D_I(\hat t)\le\theta\), where
\(\theta\) is a sufficiently small dimensional constant.  After choosing
the same center supplied by the Plan 2/FJ slicing package, let
\[
   K_{\rm in}
\]
be the union of inward fjords, holes, and same-ray missing intervals
that must be filled to obtain a star-shaped set inside a fixed annulus.
Then
\[
   |K_{\rm in}|\le C\delta .
   \tag{4.1}
\]
\end{conjecture}

\begin{conjecture}[Planar outward pruning lemma]
\label{conj:outprune}
Under the same hypotheses, let \(K_{\rm out}\subset E\) be the outward
debris that prevents the radial hull from having small \(L^\infty\)
amplitude: remote small components, active thin outward fingers, and
same-ray excess intervals beyond the annular graph radius.  Then one
can choose \(K_{\rm out}\) so that
\[
   |K_{\rm out}|\le C\delta,
   \qquad
   T(E)-T(E\setminus K_{\rm out})\le C\delta .
   \tag{4.2}
\]
\end{conjecture}

Assuming Conjectures~\ref{conj:corefill} and~\ref{conj:outprune}, the
proof of Conjecture~\ref{conj:rad} is short.  Set
\[
   E':=E\setminus K_{\rm out},\qquad
   H:=E'\cup K_{\rm in},
\]
where \(K_{\rm in}\) is chosen so that \(H\) is a small-amplitude radial
graph hull.  Then
\[
   |E\Delta H|\le |K_{\rm out}|+|K_{\rm in}|\le C\delta .
\]
By Conjecture~\ref{conj:outprune},
\[
   \dT(E')\le C\dT(E),
\]
and by monotonicity plus Lemma~\ref{lem:fillscale}, the area-normalized
copy \(G\) of \(H\) satisfies
\[
   \dT(G)\le C(\dT(E')+|K_{\rm in}|)\le C\delta .
\]
This gives (R).

\section{How far the repo gets toward the two conjectures}

\subsection{Inward defects}

Lemma~\ref{lem:caploss} is a real proof ingredient for
Conjecture~\ref{conj:corefill}.  If an inward defect \(K\) lies where
the filled torsion function is bounded below by \(a\), then it costs at
least \(a^2\capr(K)\).  In a disk, a fixed-depth slit has a fixed
positive relative capacity; a tiny hole has capacity comparable to
\(1/|\log r|\), much larger than its area.  Therefore small torsion
deficit should force the total area of fillable inward defects to be
\(O(\delta)\), often much smaller.

The missing point is not the capacity computation.  It is the
\emph{core lower bound}: one must know that the missing intervals to be
filled lie at positive torsion height in the filled comparison set.
Equivalently, one needs a coarse annular/core theorem before applying
capacity.  The current repo has measure and \(L^1\)-radial trace
control, but not this pointwise core statement.

\subsection{Outward defects}

The model outward finger estimate is favorable in dimension two.  A
thin tube of width \(w\) and length \(L\) has active-level
\[
   D_H\simeq L/w .
\]
Thus \(D_H\le\theta\) forces \(L\lesssim\theta w\), and the tube volume
\(wL\) and clearing layer \(w^2\) are of the same controlled order.
This is exactly the mechanism recorded in the Plan 3 tentacle notes.

What is missing is the arbitrary-geometry version: a theorem that
decomposes every outward non-graph part of an analytic torsion level
into such tube/finger pieces, charges them either to \(D_H\) or to
\(\dT(E)\), and proves the torsion loss after pruning is \(O(\delta)\).
Plan 2's bad-ray multiplicity estimate controls the angular measure of
extra crossings once an annular bracket is known, but it does not by
itself prove the pruning estimate (4.2).

\section{Failed shortcuts}

\subsection{Metric \(L^1\) closeness is insufficient}

The Plan 2 endpoint metric trace and the oriented radial trace control
are useful but insufficient for graph entry.  They allow a small set of
bad rays.  A literal \(L^\infty\) radial graph requires removing or
filling all bad rays, and the cost of doing so must be \(O(\delta)\), not
just \(o(1)\).

\subsection{Small isoperimetric deficit is insufficient at the needed rate}

Small \(D_I\) gives the right qualitative picture and Plan 2's bad-ray
absorption, but the Step-1 extraction at an \(O(\delta)\) layer only
gives \(D_I\le\theta\) for fixed \(\theta\), not \(D_I=O(\delta)\).
Thus \(D_I\) alone can put the level into a fixed thin annulus, but it
cannot pay for the \(O(\delta)\)-volume repair.

\subsection{The unweighted Hessian route remains false}

The shortcut
\[
   \dT(D)\ge c\int_D |D^2u_D+I/2|^2
\]
would imply all of the needed compactness at once, but it is false on
high-frequency smooth near-balls.  It must not be used.

\section{Verdict}

I cannot close Conjecture~\ref{conj:rad} from the existing repo
ingredients.  The proof has been reduced to two concrete lemmas:

\begin{enumerate}
\item a planar core/capacity fill lemma for inward fjords and holes;
\item a planar outward pruning lemma for non-graph outward debris.
\end{enumerate}

The capacity identity, Lemma~\ref{lem:caploss}, and the fill-rescale
lemma, Lemma~\ref{lem:fillscale}, are rigorous and should be retained.
They show that the intended radialization is not the wrong target.
However, the repo does not yet contain the geometric decomposition
needed to apply them to an arbitrary Step-1 torsion level.

\end{document}
